\section{Proposed System}
This chapter describes the proposed system and its feasibility. The proposed approach replaces remote PDF processing with local-first processing in the browser.

\subsection{Scope and Objective}
The scope is a browser-based PDF toolkit for routine tasks, including merge, split, organize, crop, metadata inspection, and image conversion workflows.

The main objectives are:
\begin{itemize}
  \item process PDF files without uploading documents to application servers,
  \item provide practical performance for common operations,
  \item keep the workflow simple for non-technical users,
  \item maintain transparency through an open-source codebase.
\end{itemize}

\subsection{Advantages}
Key advantages of the proposed local-first model are:
\begin{itemize}
  \item improved privacy by reducing third-party file exposure,
  \item lower dependency on network speed for processing,
  \item predictable operation in unstable connectivity scenarios,
  \item transparent behavior through inspectable source code.
\end{itemize}

\subsection{Feasibility Study}
\subsubsection{Technical Feasibility}
Modern browser APIs and mature JavaScript PDF libraries make in-browser processing feasible for a broad set of document tasks. The current project stack already supports modular feature delivery and automated testing.

\subsubsection{Economical Feasibility}
A local-first model can reduce server processing and storage overhead for document workflows. This lowers recurring infrastructure requirements compared to server-heavy processing pipelines.

\subsubsection{Operational Feasibility}
The proposed workflow matches common user expectations: select files, configure options, run action, and download output. Since processing occurs in the browser, usage is straightforward for individual users and small teams.
